% Printed172; principalPDF179.
\par\clearpage
\noindent172\hfill{\small\textbf{NACHLASS.}}\hfill\null\par\medskip
\[
\begin{aligned}
\frac{x}{a}={}&\varphi\left\{1-\frac1{2.2}ee-\frac{1.3}{2.2.4.4}e^4-\frac{1.3.3.5}{2.2.4.4.6.6}e^6-\ldots\right\}\\
&-\frac{1.3}{2.4}.ee\sin2\varphi\left\{1+\frac{1.3}{2.6}ee+\frac{1.3.3.5}{2.4.6.8}e^4+\frac{1.3.5.3.5.7}{2.4.6.6.8.10}e^6+\ldots\right\}\\
&+\frac{1.3.3.5}{2.4.6.8}.\tfrac12e^4\sin4\varphi\left\{1+\frac{3.5}{2.10}ee+\frac{3.3.5.5.7}{2.4.10.12}e^4+\frac{3.5.7.5.7.9}{2.4.6.10.12.14}e^6+\ldots\right\}\\
&-\frac{1.3.3.5.5.7}{2.4.6.8.10.12}.\tfrac13e^6\sin6\varphi\left\{1+\frac{5.7}{2.14}ee+\frac{5.7.7.9}{2.4.14.16}e^4+\frac{5.7.9.7.9.11}{2.4.6.14.16.18}e^6+\ldots\right\}
\end{aligned}
\]
{\centering etc.\par}
[oder]
\[
\begin{aligned}
&\frac{x}{a}\left\{1+\tfrac14ee+\tfrac7{64}e^4+\tfrac{15}{256}e^6+\tfrac{579}{16384}e^8+\ldots\right\}\\
&\qquad=\varphi-ee\sin2\varphi\left\{\tfrac38+\tfrac3{16}ee+\tfrac{111}{1024}e^4+[\tfrac{141}{2048}e^6+]\ldots\right\}\\
&\qquad\quad+e^4\sin4\varphi\left\{\tfrac{15}{256}+\tfrac{15}{256}ee+[\tfrac{405}{8192}e^4+]\ldots\right\}\\
&\qquad\quad-e^6\sin6\varphi\left\{\tfrac{35}{3072}+[\tfrac{35}{2048}ee+]\ldots\right\}\\
&\qquad\quad[+e^8\sin8\varphi\left\{\tfrac{315}{131072}+\ldots\right\}
\end{aligned}
\]
{\centering etc.]\par}
[Es sei nun \(\psi\) die Breite auf einer Kugel, deren Meridianquadrant dieselbe Länge wie der Meridianquadrant des Sphäroids hat, nämlich \(10\,000\,000\) Meter. Die Meridianbögen zwischen dem Äquator und den Breiten \(\varphi\) und \(\psi\) auf dem Sphäroid und auf der Kugel sollen ferner von gleicher Länge sein. Alsdann ist
\[
\psi=\frac{x\pi}{20\,000\,000}=\frac{x}{a}\left(1+\tfrac14ee+\tfrac7{64}e^4+\tfrac{15}{256}e^6+\ldots\right).
\]
Mit dem Abplattungswert \(\frac1{302{,}68}\), \(\log ee=7{,}819\,3287\), ergibt sich daher:]
\[
\begin{array}{@{}r@{\;}r@{\;}l@{\quad}l@{\;}r@{}}
\psi=\varphi-&511''{,}94191&\sin2\varphi&\log\,.\,.\,.\,.&2{,}709\,2207\\
+&0{,}52942&\sin4\varphi&&9{,}723\,80\\
-&0{,}00068&\sin6\varphi&&6{,}833\,33
\end{array}
\]
{\centering etc.\par}
Für \(\varphi\) [Göttingen] \(=51^\circ\,31'\,48''{,}7\) wird [hiernach] \(\psi=51^\circ\,23'\,29''{,}768\,245\).
[Benutzt man den Abplattungswert \(\frac1{302{,}78}\), \(\log ee=7{,}819\,1850\), so erhält man:]
\[
\begin{array}{@{}r@{\;}r@{\;}l@{\quad}l@{\;}r@{}}
\psi=\varphi-&511''{,}772\,0192&\sin2\varphi&\log\,.\,.\,.\,.&2{,}70907\,65375\\
+&0{,}529\,0745&\sin4\varphi&&9{,}72351\,68\\
-&0{,}000\,6807&\sin6\varphi&&6{,}83295\\
+&0{,}000\,0009&\sin8\varphi&&3{,}97197
\end{array}
\]
{\centering etc.\par}

% Printed173; principalPDF180.
\par\clearpage
\noindent\null\hfill{\small\textbf{KONFORME ABBILDUNG DES SPHÄROIDS IN DER EBENE.}}\hfill173\par\medskip
{\centering [22.]\par}
Wenn
\[\alpha=\beta-A\sin2\beta-B\sin4\beta-C\sin6\beta-D\sin8\beta\ldots\]
\noindent ist [worin \(A\) von der Ordnung \(ee\), \(B\) von der Ordnung \(e^4\), \(C\) von der Ordnung \(e^6\), \(D\) von der Ordnung \(e^8\), so ist bis einschließlich der achten Ordnung:]
\[
\begin{aligned}
\beta={}&\alpha+\sin2\alpha(A-AB-\tfrac12A^3)\\
&+\sin4\alpha(B+AA-2AC-4AAB-\tfrac43A^4)\\
&+\sin6\alpha(C+3AB+\tfrac32A^3)\\
&+\sin8\alpha(D+2BB+4AC+8AAB+\tfrac83A^4).
\end{aligned}
\]
[Die Umkehrung der Reihe]
\[
\begin{aligned}
\psi={}&\varphi-\sin2\varphi(\tfrac38ee+\tfrac3{16}e^4+\tfrac{111}{1024}e^6+\tfrac{141}{2048}e^8\ldots)\\
&+\sin4\varphi(\tfrac{15}{256}e^4+\tfrac{15}{256}e^6+\tfrac{405}{8192}e^8\ldots)\\
&-\sin6\varphi(\tfrac{35}{3072}e^6+\tfrac{35}{2048}e^8\ldots)\\
&+\sin8\varphi(\tfrac{315}{131072}e^8\ldots)-\text{etc.}
\end{aligned}
\]
[ist hiernach]
\[
\begin{aligned}
\varphi={}&\psi+\sin2\psi(\tfrac38ee+\tfrac3{16}e^4+\tfrac{213}{2048}e^6+\tfrac{255}{4096}e^8\ldots)\\
&+\sin4\psi(\tfrac{21}{256}e^4+\tfrac{21}{256}e^6+\tfrac{533}{8192}e^8\ldots)\\
&+\sin6\psi(\tfrac{151}{6144}e^6+\tfrac{151}{4096}e^8\ldots)\\
&+\sin8\psi(\tfrac{1097}{131072}e^8\ldots)+\text{etc.}
\end{aligned}
\]
[Da der Meridianquadrant \(=10\,000\,000\) Meter sein soll, so ist \(1\) Meter des Kugelmeridians \(=\frac{60.60.90}{10\,000\,000}\) Bogensekunden \(=0''{,}0324\). Entspricht \(\psi_0\) dem Anfangspunkt der Koordinaten, so entspricht mithin der Abszisse \(x\), wenn \(x\) nach Süden positiv ist, die geografische Breite
\[\psi=\psi_0-0''{,}0324\,x.\]
Für die Abplattung \(\frac1{302{,}68}\) ist:]
\[
\begin{array}{@{}r@{\;}r@{\;}l@{\quad}l@{\;}r@{}}
\varphi=\psi+&511''{,}94157&\sin2\psi&\log\,.\,.\,.&2{,}709\,2204\\
+&0{,}74092&\sin4\psi&&9{,}869\,77\\
+&0{,}00147&\sin6\psi&&7{,}167\,21
\end{array}
\]
{\centering etc.\par}

% Printed174; principalPDF181.
\par\clearpage
\noindent174\hfill{\small\textbf{NACHLASS.}}\hfill\null\par\medskip
{\centering [23.]\par}
Es finde zwischen der wahren Breite \(\varphi\) und der fingierten \(\omega\) folgende Relation statt
\[
\frac{\mathrm d\omega}{\cos\omega}=\frac{\mathrm d\varphi}{\cos\varphi}.\frac{1-ee}{1-ee\sin\varphi^2},
\]
\noindent woraus,
\[\sin\varphi=x,\qquad\sin\omega=y\]
\noindent gesetzt, folgt:
\[
\frac{\mathrm dy}{(1-ee)(1-yy)}=\frac{\mathrm dx}{(1-eexx)(1-xx)},
\]
\noindent also, wenn \(x\) und \(y\) zugleich verschwinden sollen,
\[
\begin{aligned}
&y+\tfrac13y^3+\tfrac15y^5+\tfrac17y^7+\tfrac19y^9\ldots\\
&\qquad=(1-ee)x+\tfrac13(1-e^4)x^3+\tfrac15(1-e^6)x^5+\tfrac17(1-e^8)x^7+\tfrac19(1-e^{10})x^9\ldots\\
&\qquad=z.
\end{aligned}
\]
Hieraus [ergibt sich] durch Umkehrung
\[
\begin{aligned}
x={}&\frac{z}{1-ee}-\frac{1+ee}{3}\left(\frac{z}{1-ee}\right)^3+\frac{2+7ee+2e^4}{15}\left(\frac{z}{1-ee}\right)^5\\
&-\frac{17+129ee+129e^4+17e^6}{315}\left(\frac{z}{1-ee}\right)^7\\
&+\frac{62+824ee+1776e^4+824e^6+62e^8}{2835}\left(\frac{z}{1-ee}\right)^9
\end{aligned}
\]
{\centering etc.\par}
Durch Substitution des Wertes von \(z\) folgt hieraus
\[
\begin{aligned}
x={}&\frac{y}{1-ee}-\frac{(3-ee)ee}{3}\left(\frac{y}{1-ee}\right)^3+\frac{(25-17ee+3e^4)e^4}{15}\left(\frac{y}{1-ee}\right)^5\\
&-\frac{(1008-1039ee+368e^4-45e^6)e^6}{315}\left(\frac{y}{1-ee}\right)^7\\
&+\frac{(18621-25758ee+13708e^4-3338e^6+315e^8)e^8}{2835}\left(\frac{y}{1-ee}\right)^9
\end{aligned}
\]
{\centering etc.\par}
Setzt man
\[\frac{ee}{1-ee}=\delta,\]

% Printed175; principalPDF182.
\par\clearpage
\noindent\null\hfill{\small\textbf{KONFORME ABBILDUNG DES SPHÄROIDS IN DER EBENE.}}\hfill175\par\medskip
\noindent so verwandelt sich dieser Ausdruck in folgenden:
\[
\begin{aligned}
x=(1+\delta)\biggl\{y&-\frac{3\delta+2\delta\delta}{3}y^3+\frac{25\delta\delta+33\delta^3+11\delta^4}{15}y^5\\
&-\frac{1008\delta^3+1985\delta^4+1314\delta^5+292\delta^6}{315}y^7\\
&+\frac{18621\delta^4+48726\delta^5+48160\delta^6+21288\delta^7+3548\delta^8}{2835}y^9\text{ etc.}\biggr\}.
\end{aligned}
\]
Hieraus folgt ferner
\[
\begin{aligned}
xx=(1+\delta)^2\biggl\{yy&-\frac{6\delta+4\delta\delta}{3}y^4+\frac{195\delta\delta+258\delta^3+86\delta^4}{45}y^6\\
&-\frac{3066\delta^3+6056\delta^4+4014\delta^5+892\delta^6}{315}y^8\\
&+\frac{316305\delta^4+830340\delta^5+822217\delta^6+363738\delta^7+60623\delta^8}{14175}y^{10}\text{ etc.}\biggr\}
\end{aligned}
\]
[und, indem man den daraus sich ergebenden Ausdruck für \(1-xx\) noch mit \((1-yy)(1+yy+y^4+y^6+\ldots)\) multipliziert,]
\[
\begin{aligned}
1-xx=(1-yy)\biggl\{1&-(2\delta+\delta\delta)yy+\frac{13\delta\delta+14\delta^3+4\delta^4}{3}y^4\\
&-\frac{438\delta^3+737\delta^4+430\delta^5+86\delta^6}{45}y^6\\
&+\frac{7029\delta^4+16182\delta^5+14374\delta^6+5798\delta^7+892\delta^8}{315}y^8\text{ etc.}\biggr\}.
\end{aligned}
\]
Hieraus ferner
\[
\begin{aligned}
\sqrt{\frac{1-xx}{1-yy}}=\frac{\cos\varphi}{\cos\omega}={}&1-\frac{2\delta+\delta\delta}{2}yy+\frac{40\delta\delta+44\delta^3+13\delta^4}{24}y^4\\
&-\frac{2304\delta^3+3976\delta^4+2390\delta^5+493\delta^6}{720}y^6
\end{aligned}
\]
{\centering etc.\par}
Da nun
\[
\begin{aligned}
\frac{1-ee\sin\varphi^2}{1-ee}=1+\delta\cos\varphi^2=[(1+\delta)-\delta xx={}]\,(1+\delta)\biggl\{1&-(\delta+\delta\delta)yy+\frac{6\delta\delta+10\delta^3+4\delta^4}{3}y^4\\
&-\frac{195\delta^3+453\delta^4+344\delta^5+86\delta^6}{45}y^6\text{ etc.}\biggr\},
\end{aligned}
\]
\noindent so wird
\[
\begin{aligned}
\frac{\mathrm d\varphi}{\mathrm d\omega}=(1+\delta)\biggl\{1&-\frac{4\delta+3\delta\delta}{2}yy+\frac{112\delta\delta+160\delta^3+57\delta^4}{24}y^4\\
&-\frac{8064\delta^3+16864\delta^4+11764\delta^5+2739\delta^6}{720}y^6\text{ etc.}\biggr\}.
\end{aligned}
\]
Hiernach ergibt die Integration
\[ [\varphi=\omega+c_1\sin\omega\cos\omega+c_2\sin\omega^3\cos\omega+c_3\sin\omega^5\cos\omega\ldots,\]

% Printed176; principalPDF183.
\par\clearpage
\noindent176\hfill{\small\textbf{NACHLASS.}}\hfill\null\par\medskip
\noindent also wird,] da
\[
\begin{gathered}
\mathrm d(\sin\omega\cos\omega)=(1-2yy)\mathrm d\omega,\qquad
\mathrm d(\sin\omega^3\cos\omega)=(3yy-4y^4)\mathrm d\omega,\\
\mathrm d(\sin\omega^5\cos\omega)=(5y^4-6y^6)\mathrm d\omega,\text{ etc.}
\end{gathered}
\]
\noindent ist,
\[
[\frac{\mathrm d\varphi}{\mathrm d\omega}=1+c_1+(-2c_1+3c_2)yy+(-4c_2+5c_3)y^3+(-6c_3+7c_4)y^4\ldots.
\]
Der Vergleich mit dem vorigen Ausdruck für \(\frac{\mathrm d\varphi}{\mathrm d\omega}\) liefert die Werte der Koeffizienten \(c_1,c_2,c_3,\ldots\) und damit]
\[
\begin{aligned}
\varphi={}&\omega+\delta\sin\omega\cos\omega-\frac{7\delta\delta+3\delta^3}{6}\sin\omega^3\cos\omega\\
&+\frac{224\delta^3+217\delta^4+57\delta^5}{120}\sin\omega^5\cos\omega\\
&-\frac{17116\delta^4+26576\delta^5+14503\delta^6+2739\delta^7}{5040}\sin\omega^7\cos\omega
\end{aligned}
\]
{\centering etc.\par}
\bigskip
{\centering [24.]\par}
Leichter ist die umgekehrte Aufgabe, [\(\omega\) aus der Gleichung
\[
\frac{\mathrm d\omega}{\mathrm d\varphi}=\sqrt{\frac{1-yy}{1-xx}}.\frac{1-ee}{1-eexx}
\]
\noindent abzuleiten,] wobei sich Folgendes ergibt:
\[
\begin{aligned}
\omega={}&\varphi-ee\sin\varphi\cos\varphi-\frac{5e^4-e^6}{6}\sin\varphi^3\cos\varphi\\
&-\frac{104e^6-45e^8+5e^{10}}{120}\sin\varphi^5\cos\varphi
-\frac{4948e^8[-3360e^{10}+777e^{12}-61e^{14}]}{5040}\sin\varphi^7\cos\varphi
\end{aligned}
\]
{\centering etc.\par}
Ferner wird, wenn \(\psi\) dem Bogen proportional ist, dieser \(=K\psi\), [wobei]
\[K=a(1-\tfrac14ee-\tfrac3{64}e^4-\tfrac5{256}e^6-\tfrac{175}{16384}e^8\ldots),\]
[und nach Art.~22]
\[
\begin{aligned}
\psi={}&\varphi-\sin\varphi\cos\varphi(\tfrac34ee+\tfrac9{64}e^4+\tfrac{13}{256}e^6+\tfrac{381}{16384}e^8\ldots)\\
&-\sin\varphi^3\cos\varphi(\tfrac{15}{32}e^4+\tfrac5{48}e^6+\tfrac{335}{8192}e^8\ldots)\\
&-\sin\varphi^5\cos\varphi(\tfrac{35}{96}e^6+\tfrac{175}{2048}e^8\ldots)\\
&-\sin\varphi^7\cos\varphi(\tfrac{315}{1024}e^8\ldots)
\end{aligned}
\]
{\centering etc.;\par}
[folglich hat man zunächst:]

% Printed177; principalPDF184.
\par\clearpage
\noindent\null\hfill{\small\textbf{KONFORME ABBILDUNG DES SPHÄROIDS IN DER EBENE.}}\hfill177\par\medskip
\[
\begin{aligned}
\psi={}&\omega+\sin\varphi\cos\varphi(\tfrac14ee-\tfrac9{64}e^4-\tfrac{15}{256}e^6-\tfrac{381}{16384}e^8\ldots)\\
&+\sin\varphi^3\cos\varphi(\tfrac{35}{96}e^4-\tfrac{13}{48}e^6-\tfrac{335}{8192}e^8\ldots)\\
&+\sin\varphi^5\cos\varphi(\tfrac{211}{480}e^6-\tfrac{913}{2048}e^8\ldots)\\
&+\sin\varphi^7\cos\varphi(\tfrac{217447}{322560}e^8\ldots)
\end{aligned}
\]
{\centering etc.\par}
Hier genügt es, [nach S.~174 und 175] zu setzen:
\[
\begin{aligned}
\mathopen{[}x={}]\sin\varphi={}&\sin\omega\left\{(1+ee+e^4+e^6)-\frac{3ee+8e^4+15e^6}{3}yy+\frac{25e^4+108e^6}{15}y^4-\frac{16e^6}{5}y^6\right\}\\
\cos\varphi={}&\cos\omega\left\{1-\frac{2ee+3e^4+4e^6}{2}yy+\frac{10e^4+31e^6}{6}y^4-\frac{16e^6}{5}y^6\right\},
\end{aligned}
\]
[also]
\[
\begin{aligned}
\sin\varphi\cos\varphi=\sin\omega\cos\omega\biggl\{1+ee+e^4+e^6&-\frac{12ee+31e^4+57e^6}{6}yy\\
&+\frac{65e^4+273e^6}{15}y^4-\frac{146}{15}e^6y^6\biggr\}
\end{aligned}
\]
\[
\begin{aligned}
\sin\varphi^3\cos\varphi&=\sin\omega^3\cos\omega\biggl\{1+3ee+6e^4-\biggl[\frac{8ee+37e^4}{2}yy+\frac{38}{3}e^4y^4\biggr\}\\
\sin\varphi^5\cos\varphi&=\sin\omega^5\cos\omega\{1+5ee-6eeyy\}\\
\sin\varphi^7\cos\varphi&=\sin\omega^7\cos\omega.
\end{aligned}
\]
Damit wird:
\[
\begin{aligned}
\psi={}&\omega+\sin\omega\cos\omega(\tfrac14ee+\tfrac7{64}e^4+\tfrac{15}{256}e^6+\tfrac{579}{16384}e^8\ldots)\\
&-\sin\omega^3\cos\omega(\tfrac{13}{96}e^4+\tfrac3{16}e^6+\tfrac{1743}{8192}e^8\ldots)\\
&+\sin\omega^5\cos\omega(\tfrac{61}{480}e^6+\tfrac{10111}{30720}e^8\ldots)\\
&-\sin\omega^7\cos\omega(\tfrac{49561}{322560}e^8\ldots)
\end{aligned}
\]
{\centering etc.\par}
oder]
\[
\begin{aligned}
\psi={}&\omega+\sin2\omega(\tfrac18ee+\tfrac1{48}e^4+\tfrac7{3072}e^6-\tfrac{83}{92160}e^8\ldots)\\
&+\sin4\omega(\tfrac{13}{768}e^4+\tfrac{29}{3840}e^6+\tfrac{833}{368640}e^8\ldots)\\
&+\sin6\omega(\tfrac{61}{15360}e^6+\tfrac{221}{71680}e^8\ldots)\\
&+\sin8\omega(\tfrac{49561}{41287680}e^8\ldots)
\end{aligned}
\]
{\centering etc.\par}
[Aus der Schlussformel des Art.~23 erhält man]
\par\noindent{\small IX.}\hfill{\small23}\hspace*{.12\linewidth}\par

% Printed178; principalPDF185.
\par\clearpage
\noindent178\hfill{\small\textbf{NACHLASS.}}\hfill\null\par\medskip
\[
\begin{aligned}
\varphi={}&\omega+\sin2\omega(\tfrac12\delta-\tfrac7{24}\delta\delta+\tfrac16\delta^3-\tfrac4{45}\delta^4\ldots)\\
&+\sin4\omega(\tfrac7{48}\delta\delta-\tfrac{41}{240}\delta^3+\tfrac{335}{2304}\delta^4\ldots)\\
&+\sin6\omega(\tfrac7{120}\delta^3-\tfrac{23}{224}\delta^4\ldots)\\
&+\sin8\omega(\tfrac{4279}{161280}\delta^4\ldots)
\end{aligned}
\]
{\centering etc.\par}
\[\delta=\frac{ee}{1-ee}\]
[und aus der ersten Formel dieses Artikels:]
\[
\begin{aligned}
\omega={}&\varphi-\sin2\varphi(\tfrac12ee+\tfrac5{24}e^4+\tfrac3{32}e^6+\tfrac{281}{5760}e^8\ldots)\\
&+\sin4\varphi(\tfrac5{48}e^4+\tfrac7{80}e^6+\tfrac{697}{11520}e^8\ldots)\\
&-\sin6\varphi(\tfrac{13}{480}e^6+\tfrac{461}{13440}e^8\ldots)\\
&+\sin8\varphi(\tfrac{1237}{161280}e^8\ldots)
\end{aligned}
\]
{\centering etc.\par}
\bigskip
{\centering [25.]\par}
[Setzt man
\[\tfrac14e=f,\]
\noindent also
\[f=\frac{\sqrt{2n-1}}{4n},\]
\noindent wobei \(\frac1n\) die Abplattung bedeutet, so wird [nach Art.~22 und 24]:]
\[
\begin{aligned}
\psi={}&\varphi-(6ff+48f^4+444f^6+4512f^8\ldots)\sin2\varphi\\
&+(15f^4+240f^6+3240f^8\ldots)\sin4\varphi\\
&-(\tfrac{140}{3}f^6+1120f^8\ldots)\sin6\varphi\\
&+(\tfrac{315}{2}f^8\ldots)\sin8\varphi
\end{aligned}
\]
{\centering etc.\par}
[Nach S.~172 oder 176 ist \(a(1-\tfrac14ee-\tfrac3{64}e^4\ldots)\mathrm d\psi=\mathrm dx\), daher wird]
\[
\frac{\mathrm d\psi}{\mathrm d\varphi}(1-4ff-12f^4-80f^6-700f^8\ldots)=\frac{1-16ff}{(1-16ff\sin\varphi^2)^{\frac32}},
\]
[Weiter ist]

% Printed179; principalPDF186.
\par\clearpage
\noindent\null\hfill{\small\textbf{KONFORME ABBILDUNG DES SPHÄROIDS IN DER EBENE.}}\hfill179\par\medskip
\[
\begin{aligned}
\varphi={}&\psi+(6ff+48f^4+426f^6+4080f^8\ldots)\sin2\psi\\
&+(21f^4+336f^6+4264f^8\ldots)\sin4\psi\\
&+(\tfrac{302}{3}f^6+2416f^8\ldots)\sin6\psi\\
&+(\tfrac{1097}{2}f^8\ldots)\sin8\psi
\end{aligned}
\]
{\centering etc.\par}
\[
\begin{aligned}
\varphi={}&\omega+(8ff+\tfrac{160}{3}f^4+\tfrac{1024}{3}f^6+\tfrac{106496}{45}f^8\ldots)\sin2\omega\\
&+(\tfrac{112}{3}f^4+\tfrac{7424}{15}f^6+\tfrac{207616}{45}f^8\ldots)\sin4\omega\\
&+(\tfrac{3584}{15}f^6+\tfrac{165888}{35}f^8\ldots)\sin6\omega\\
&+(\tfrac{547712}{315}f^8\ldots)\sin8\omega
\end{aligned}
\]
{\centering etc.\par}
\[
\begin{aligned}
\omega={}&\varphi-(8ff+\tfrac{160}{3}f^4+384f^6+\tfrac{143872}{45}f^8\ldots)\sin2\varphi\\
&+(\tfrac{80}{3}f^4+\tfrac{1792}{5}f^6+\tfrac{178432}{45}f^8\ldots)\sin4\varphi\\
&-(\tfrac{1664}{15}f^6+\tfrac{236032}{105}f^8\ldots)\sin6\varphi\\
&+(\tfrac{158336}{315}f^8\ldots)\sin8\varphi
\end{aligned}
\]
{\centering etc.\par}
\[
\begin{aligned}
\psi={}&\omega+(2ff+\tfrac{16}{3}f^4+\tfrac{28}{3}f^6-\tfrac{2656}{45}f^8\ldots)\sin2\omega\\
&+(\tfrac{13}{3}f^4+\tfrac{464}{15}f^6+\tfrac{6664}{45}f^8\ldots)\sin4\omega\\
&+(\tfrac{244}{15}f^6+\tfrac{7072}{35}f^8\ldots)\sin6\omega\\
&+(\tfrac{49561}{630}f^8\ldots)\sin8\omega
\end{aligned}
\]
{\centering etc.\par}
[und hieraus, indem man die Formel von S.~173, oben, benutzt:]
\[
\begin{aligned}
\omega={}&\psi-(2ff+\tfrac{16}{3}f^4+14f^6-\tfrac{272}{45}f^8\ldots)\sin2\psi\\
&-(\tfrac13f^4+\tfrac{48}{5}f^6+\tfrac{4472}{45}f^8\ldots)\sin4\psi\\
&-(\tfrac{34}{15}f^6+\tfrac{4528}{105}f^8\ldots)\sin6\psi\\
&-(\tfrac{4397}{630}f^8\ldots)\sin8\psi
\end{aligned}
\]
{\centering etc.\par}
\par\noindent\null\hfill{\small23*}\hspace*{.12\linewidth}\par

% Printed180; principalPDF187.
\par\clearpage
\noindent180\hfill{\small\textbf{NACHLASS.}}\hfill\null\par\medskip
{\centering [26.]\par\medskip Zur Berechnung von \(\log\cos\varphi\).\par}
Ist
\[p=q+h\sin q,\]
\noindent so wird
\[\log\sin p=\log\sin q+h\cos q-\tfrac12hh+\tfrac13h^3\cos q-\tfrac1{12}h^4(1+2\cos q^2)\ldots.\]
[Die Logarithmen sind hier wie auch in Art.~27 hyperbolisch.]

Also,
\[h=a\cos q+b\cos3q+c\cos5q+d\cos7q\ldots\]
\noindent gesetzt,
\[
\begin{aligned}
\log\sin p={}&\log\sin q\\
&+(\tfrac12a-\tfrac14aa+\tfrac18a^3-\tfrac1{12}a^4+\tfrac18aab-\tfrac14bb\ldots)\\
&+\cos2q(\tfrac12a-\tfrac14aa+\tfrac12b-\tfrac12ab+\tfrac16a^3-\tfrac{23}{192}a^4+\tfrac38aab\ldots)\\
&+\cos4q(\tfrac12b-\tfrac12ab+\tfrac12c+\tfrac1{24}a^3-\tfrac1{24}a^4+\tfrac38aab-\tfrac12ac\ldots)\\
&+\cos6q(\tfrac12c-\tfrac1{192}a^4+\tfrac18aab-\tfrac12ac-\tfrac14bb+\tfrac12d\ldots)\\
&+\cos8q(\tfrac12d\ldots)
\end{aligned}
\]
{\centering etc.\par}
In unserem Fall ist [\(p=90^\circ-\varphi\) und \(q=90^\circ-\psi\) zu setzen. Geht man dann von der zweiten Reihe zwischen \(\varphi\) und \(\psi\) in Art.~25 aus, erhält man:]
\[
\begin{aligned}
a&=-12ff-54f^4-\frac{1144}{3}f^6-3367f^8\ldots\\
b&=\phantom{-12ff-}+42f^4+\frac{1412}{3}f^6+4793f^8\ldots\\
c&=\phantom{-12ff-54f^4}-\frac{604}{3}f^6-3735f^8\ldots\\
d&=\phantom{-12ff-54f^4-\frac{1144}{3}f^6}+1097f^8\ldots
\end{aligned}
\]
{\centering etc.\par}
Also:
\[
\begin{aligned}
\log\cos\varphi={}&\log\cos\psi-6ff-63f^4-\tfrac{2192}{3}f^6-\tfrac{18059}{2}f^8\ldots\\
&+(6ff+42f^4+\tfrac{946}{3}f^6+2450f^8\ldots)\cos2\psi\\
&+(21f^4+\tfrac{944}{3}f^6+3711f^8\ldots)\cos4\psi\\
&+(\tfrac{302}{3}f^6+2320f^8\ldots)\cos6\psi\\
&+(\tfrac{1097}{2}f^8\ldots)\cos8\psi
\end{aligned}
\]
{\centering etc.\par}

% Printed181; principalPDF188.
\par\clearpage
\noindent\null\hfill{\small\textbf{KONFORME ABBILDUNG DES SPHÄROIDS IN DER EBENE.}}\hfill181\par\medskip
{\centering [27.]\par\medskip
\(\left[\text{Berechnung von }\log(1-ee\sin\varphi^2),\quad
\log\frac{(1-ee\sin\varphi^2)^2}{aa(1-ee)}\text{ und }\frac{a\cos\varphi}{\sqrt{1-ee\sin\varphi^2}}.\right]\)\par}
Man hat ferner [wenn man berücksichtigt, dass
\[
\begin{aligned}
\sin\varphi={}&(1+3ff+15f^4+87f^6+528f^8\ldots)\sin\psi\\
&+(3ff+30f^4+264f^6+2235f^8\ldots)\sin3\psi\\
&+(15f^4+\tfrac{790}{3}f^6+3360f^8\ldots)\sin5\psi\\
&+(\tfrac{259}{3}f^6+2184f^8\ldots)\sin7\psi\\
&+(531f^8\ldots)\sin9\psi
\end{aligned}
\]
{\centering etc.\par}
\noindent wird:]
\[
\begin{aligned}
\log(1-ee\sin\varphi^2)={}&-8ff-96f^4-\tfrac{3584}{3}f^6-15424f^8\ldots\\
&+(8ff+64f^4+520f^6+4288f^8\ldots)\cos2\psi\\
&+(32f^4+512f^6+\tfrac{18944}{3}f^8\ldots)\cos4\psi\\
&+(\tfrac{488}{3}f^6+3904f^8\ldots)\cos6\psi\\
&+(\tfrac{2752}{3}f^8\ldots)\cos8\psi
\end{aligned}
\]
{\centering etc.\par}
\noindent und für den Logarithmus des Krümmungsmaßes:
\[
\begin{aligned}
-2\log a\qquad\ast\quad&-64f^4-1024f^6-14464f^8\ldots\\
&+(16ff+128f^4+1040f^6+8576f^8\ldots)\cos2\psi\\
&+(64f^4+1024f^6+\tfrac{37888}{3}f^8\ldots)\cos4\psi\\
&+(\tfrac{976}{3}f^6+7808f^8\ldots)\cos6\psi\\
&+(\tfrac{5504}{3}f^8\ldots)\cos8\psi
\end{aligned}
\]
{\centering etc.\par}
[Für das Krümmungsmaß selbst ergibt sich:]
\[
\begin{aligned}
\frac1{aa}.\frac{(1-16ff\sin\varphi^2)^2}{1-16ff}=\frac1{aa}\biggl\{1&+(16ff+128f^4+1040f^6+8576f^8\ldots)\cos2\psi\\
&+(128f^4+2048f^6+\tfrac{74752}{3}f^8\ldots)\cos4\psi\\
&+(1008f^6+24192f^8\ldots)\cos6\psi\\
&+(\tfrac{23552}{3}f^8\ldots)\cos8\psi\text{ etc.}\biggr\}
\end{aligned}
\]
